\section{기술적 파급효과와 향후 추가 개발 계획}

\subsection{기술적 파급효과}

\subsubsection{교육적 파급효과}

\begin{itemize}
  \item \textbf{진입 장벽 완화}: 시각적 블록과 한국어 UI, 어린이용 단순화
        인터페이스로 초등 저학년도 물리 로봇 제어를 경험한다.
  \item \textbf{실기 없는 학습 보장}: 동일 블록이 동일 명령으로 시뮬레이터와
        실기 양쪽에서 동작하므로, 로버 보급 대수에 구애받지 않고 수업을
        진행할 수 있다.
  \item \textbf{단계적 개념 습득}: 디지털 출력$\to$다채널$\to$난수$\to$구동$\to$
        센서$\to$제어흐름$\to$모듈화로 이어지는 미션 설계가 컴퓨팅 사고력을
        점진적으로 형성한다.
\end{itemize}

\subsubsection{기술적 파급효과}

\begin{itemize}
  \item \textbf{오프라인 우선 설계}: esbuild 번들과 fetch 심을 통한 \code{file://}
        단독 실행은, 네트워크\textbf{·}계정\textbf{·}설치가 제약되는 교실 환경에
        적합한 배포 모델을 제시한다. 이는 다른 웹 기반 교구에 재사용 가능한
        패턴이다.
  \item \textbf{개방형 하드웨어 구성}: \code{system.json} 기반 모듈\textbf{·}핀
        런타임 구성은 동일 펌웨어로 다양한 하드웨어 조합을 수용하여,
        저비용\textbf{·}맞춤형 교구 제작을 가능케 한다.
  \item \textbf{표준 웹 기술 활용}: Web Bluetooth\textbf{·}WebGL\textbf{·}Blockly 등
        표준 기술만으로 구성되어 특정 벤더 종속이 없고, 유지보수\textbf{·}확장이
        용이하다.
\end{itemize}

\subsection{향후 추가 개발 계획}

\subsubsection{단기 (기존 자산 강화)}

\begin{itemize}
  \item 통신 계층 추상화 및 명령 프로토콜 단일 출처화(11장 우선순위 항목).
  \item AI 도우미(NL$\to$블록, 코딩 질의응답)의 편집기 정식 통합.
  \item 펌웨어 시간 명령 비차단화로 비상정지 반응성 개선.
  \item 단위 테스트\textbf{·}빌드 스모크 테스트 도입.
\end{itemize}

\subsubsection{중기 (플랫폼 확장)}

\begin{itemize}
  \item Electron 데스크톱 앱 출시(코드 재사용 최대).
  \item Android PWA 제공 및 Capacitor 래퍼 검토.
  \item iOS는 Capacitor + 네이티브 BLE 플러그인으로 대응.
  \item 미션\textbf{·}예제 콘텐츠 확장 및 교사용 진도\textbf{·}평가 도구 보강.
\end{itemize}

\subsubsection{장기 (플랫폼 고도화)}

\begin{itemize}
  \item 다중 로버\textbf{·}협동 미션(다중 BLE 연결) 지원 검토.
  \item 센서\textbf{·}구동 모듈 확장(라인트레이서, 카메라 등)과 그에 맞는 블록\textbf{·}
        명령 추가.
  \item 학습 데이터\textbf{·}작업물 클라우드 동기화(선택적, 오프라인 우선 원칙
        유지).
  \item 다국어 지원을 통한 해외 보급.
  \item 시뮬레이터의 물리\textbf{·}센서 모델 정교화로 실기와의 정합성 향상.
\end{itemize}

\subsection{개발 로드맵 개관}

위 계획을 단기\textbf{·}중기\textbf{·}장기로 배치하면 그림~\ref{fig:roadmap}와 같다.
전송 추상화와 프로토콜 단일 출처화가 단기 토대가 되고, 이를 발판으로 중기의
멀티플랫폼 확장과 장기의 플랫폼 고도화가 이어진다.

\begin{diagram}{단기 → 중기 → 장기 개발 로드맵}{fig:roadmap}
\node[dom, text width=42mm] (s) {단기\\{\footnotesize\mdseries 전송 추상화 · 프로토콜 단일화 · AI 통합 · 테스트}};
\node[dom, text width=42mm, right=10mm of s] (m) {중기\\{\footnotesize\mdseries Electron · Android PWA/Capacitor · iOS 네이티브 BLE · 콘텐츠 확장}};
\node[dom, text width=42mm, right=10mm of m] (l) {장기\\{\footnotesize\mdseries 다중 로버 · 모듈 확장 · 다국어 · 시뮬레이터 고도화}};
\draw[flow] (s) -- node[lbl,above]{토대} (m);
\draw[flow] (m) -- node[lbl,above]{확장} (l);
\node[hw, below=8mm of s, text width=42mm] (sf) {기존 자산 강화};
\node[hw, below=8mm of m, text width=42mm] (mf) {플랫폼 확장};
\node[hw, below=8mm of l, text width=42mm] (lf) {플랫폼 고도화};
\draw[dep] (s)--(sf); \draw[dep] (m)--(mf); \draw[dep] (l)--(lf);
\end{diagram}

\begin{teachbox}{이 설계가 남기는 것}
\teachhead{설계 의도}
\begin{itemize}
  \item 오프라인 우선\textbf{·}표준 웹 기술\textbf{·}개방형 하드웨어 구성이 교육 보급성과 확장성을 만든다.
\end{itemize}
\teachhead{분석할 때 핵심 질문}
\begin{itemize}
  \item 이 설계의 어떤 점이 \emph{다른 교구}에도 재사용될 수 있는가?
  \item 새 기능을 더할 때 ``오프라인 우선'' 원칙을 어떻게 지킬 것인가?
\end{itemize}
\teachhead{점검 요소}
\begin{itemize}
  \item 새 미션이나 새 하드웨어 모듈을 하나 제안하고, 1\textasciitilde 12장 중 \emph{어느 장}의 코드를 수정해야 하는지 매핑하라.
\end{itemize}
\end{teachbox}

\subsection{결론}

ARES는 시각적 블록 코딩, 표준 웹 통신(Web Bluetooth), WebGL 3D 시뮬레이션,
모듈식 MicroPython 펌웨어를 결합하여, 초등 교육 현장에 적합한 오프라인 우선
로봇 코딩 플랫폼을 구현하였다. 코드베이스는 역할별로 명확히 분리되어 있으며,
웹과 펌웨어를 잇는 단순 텍스트 프로토콜이 시스템의 유연성과 확장성의 근간을
이룬다. 본 보고서가 분석한 구조적 특성과 개선\textbf{·}확장 계획은, ARES를
네이티브 멀티플랫폼으로 확장하고 콘텐츠\textbf{·}AI 보조 기능을 고도화하는
다음 단계의 기술적 토대가 된다.
